\frontmatter												%Seitennumerierung
\pagenumbering{Roman}										%Römische Zahlen
\addtocounter{page}{2}

\newcommand{\doublesignature}[2]{%
  \parbox{\textwidth}{
    \hfill
    \parbox{7cm}{
      \centering
      \rule{6cm}{1pt}\\
      #1
    }
    \parbox{7cm}{
      \centering
      \rule{6cm}{1pt}\\
      #2
    }
  }
  \mbox{}\\
  \mbox{}\\
  \mbox{}\\
  \mbox{}\\
}

\vspace*{20pt}

\section*{Eidesstattliche Erklärung}
\label{sec:eidestattliche-erklaerung}
Ich erkläre an Eides statt, dass ich die vorliegende Arbeit selbstständig verfasst, andere als die angegebenen
Quellen/Hilfsmittel nicht benutzt und die den benutzten Quellen wörtlich und inhaltlich entnommenen
Stellen als solche kenntlich gemacht habe.\\
\\
Arnfels, am 05. April 2018\\

\vskip 2cm

\doublesignature{Lukas Freyler}{Stefan Ornik}
\doublesignature{Fabio Pölzl}{Dominik Riegelnegg}

\vskip 5cm

\clearpage

\newpage
\thispagestyle{empty}
\mbox{}

\clearpage

\section*{Danksagung}
\label{sec:danksagung}
An dieser Stelle möchten wir uns bei allen bedanken, die uns im Rahmen der Diplomarbeit unterstützt und betreut haben.\\
\\
Besonders bedanken möchten wir uns bei unseren Betreuern, Dipl.-Ing. Werner Harnisch, Dipl.-Ing. Wolfgang Mader, BEd. Otto Schuller
und Dipl.-Ing. Manfred Steiner für die fachliche Unterstützung während dieser Arbeit.

\newpage

\section*{Zusammenfassung}

Das Hauptziel der Diplomarbeit ist es, Pferden mehr Bewegung zu ermöglichen. In der freien Natur legen Pferde einige Kilometer am Tag zurück. Dies können die Pferde in den verschiedensten Reitställen nicht behaupten. Die Haltung eines Pferdes bringt viele Kosten mit sich. Da bleibt meistens nichts mehr für ein neues Trainingsgerät für die Pferde übrig. Als Nächstes kommt hinzu, dass manche Personen auch nicht die nötige Zeit haben, um neben der Pferdeführanlage zu stehen und den Betrieb zu überwachen. Dadurch ergeben sich verschiedene Punkte, um die Schwierigkeiten für den Besitzer zu minimieren: \\
Unser Ziel ist es eine automatisierte Anlage zu entwickeln, welche einfach über eine Android Applikation zu steuern ist. Da die Anlage auch an heißen Tagen betrieben wird, wurde eine Kühlung für die Pferde auf die Liste der Abarbeitungspunkte gesetzt. Die Erfrischung soll mittels Sprühregen realisiert werden, damit die Pferde bei heißen Temperaturen eine Abkühlung bekommen können. Der Benutzer kann bis zu vier Pferde innerhalb der Anlage positionieren und diese starten. Die Zeit und die Geschwindigkeit der Pferdeführanlage soll er selbst auswählen können. Als letzten Punkt soll es eine Videoübertragung zur Android Applikation geben, damit der Benutzer jederzeit einen Blick auf die Anlage werfen kann, ob noch alles in Ordnung ist.  \newline{}
 
Was sind nun die wichtigsten Punkte von unserer Anlage?\\

\begin{itemize}
\item{Automatisierte Anlage}
\item{Platz für vier Pferde gleichzeitig}
\item{Android Applikation um die Anlage zu steuern}
\item{Videoüberwachung, die Bestandteil der Applikation ist}
\item{Einfache Bedienung}
\item{Geschwindigkeit und Zeit sollen individuell konfigurierbar sein} 
\item{Sprühregen zur Abkühlung der Pferde}
\item{Langlebig}
\end{itemize}

\newpage

\section*{Abstract}
\label{sec:abstractIntroduction}

The primary goal of this thesis is to enable horses more movement. Horses travel many kilometers a day in the case they live in the wildlife. Horses which are the most time of their life inside the stable are not able to do so. Moreover, keeping those animals implicates a lot of costs. As a result there is not any money left for a new training device. Another difficult is, that the keepers of the animals do not have the time required to stand by the horse exerciser and control the system. \\

This thesis wants to find ways to minimize those problems. The aim of the thesis is to develop an automated system which can be easily controlled by an Android application. The program should allow the user to compile a program for the horses. Because of the reason, that the horse exerciser is also going to operate during hot days, one mayor point is to develop a cooling system for the horses. If the user wants to train some horses, he has to place up to four horses inside the horse exerciser. Then he is able to start the system. He should also be able to set the operating time and speed of the horse exerciser. As a final step, as the system is video controlled, the user can check on the app, if everything is okay. \newline{}

What are the primary goals of the horse exerciser?\\

\begin{itemize}
	\item{automated unit}
	\item{capacity for up to four horses simultaneously}
	\item{Android application to control the system}
	\item{video control which is available inside the app}
	\item{easy handling}
	\item{speed and time should be able to be adapted individually}
	\item{light drizzle to cool down the horses}
	\item{persistent}
\end{itemize} 


\clearpage

\newpage
\thispagestyle{empty}
\mbox{}

\clearpage

\subsection*{Gender Erklärung}
\label{sec:gender-erklaerung}
Aus Gründen der besseren Lesbarkeit wird in dieser Arbeit die Sprachform des generischen Maskulinums angewendet. Es wird an dieser Stelle darauf hingewiesen, dass die ausschließliche Verwendung der männlichen Form geschlechtsunabhängig verstanden werden soll.

\subsection*{Über dieses Dokument}
\label{sec:ueber-dokument}
Diese Arbeit wurde in \LaTeX{} verfasst. Diese Art der Dokumentation bietet gegenüber den normalen Textverarbeitungen gewisse Vorteile hinsichtlich der Formatierung und des Einbindens von Grafiken. Auch Formeln können sehr einfach und effizient angegeben werden.

\clearpage

\newpage
\thispagestyle{empty}
\mbox{}

\clearpage

\section*{Projektteam}
\label{sec:projektteam}

\subsection*{Lukas Freyler}
\begin{wrapfigure}[10]{l}{0.5\textwidth}
\begin{center}
  \includegraphics[width=0.35\textwidth]{fig/logoMecha}
\end{center}
\end{wrapfigure}
\mbox{}\\
\mbox{}\\
\textbf{Aufgabenbereich}:\\
\LaTeX{}\\
\textbf{Betreuer}:\\
Dipl.-Ing. Manfred Steiner
\mbox{}\\
\mbox{}\\
\mbox{}\\
\mbox{}\\
\mbox{}\\
\mbox{}\\

\subsection*{Stefan Ornik}
\begin{wrapfigure}[10]{l}{0.5\textwidth}
\begin{center}
  \includegraphics[width=0.35\textwidth]{fig/logoMecha}
\end{center}
\end{wrapfigure}
\mbox{}\\
\mbox{}\\
\textbf{Aufgabenbereich}:\\
\LaTeX{}\\
\textbf{Betreuer}:\\
Dipl.-Ing. Wolfgang Mader
\mbox{}\\
\mbox{}\\
\mbox{}\\
\mbox{}\\
\mbox{}\\
\newpage

\subsection*{Fabio Pölzl}
\begin{wrapfigure}[10]{l}{0.5\textwidth}
\begin{center}
  \includegraphics[width=0.35\textwidth]{fig/logoMecha}
\end{center}
\end{wrapfigure}
\mbox{}\\
\mbox{}\\
\textbf{Aufgabenbereich}:\\
\LaTeX{}\\
\textbf{Betreuer}:\\
Dipl.-Ing. Werner Harnisch
\mbox{}\\
\mbox{}\\
\mbox{}\\
\mbox{}\\
\mbox{}\\
\mbox{}\\

\subsection*{Dominik Riegelnegg}
\begin{wrapfigure}[10]{l}{0.5\textwidth}
\begin{center}
  \includegraphics[width=0.35\textwidth]{fig/logoMecha}
\end{center}
\end{wrapfigure}
\mbox{}\\
\mbox{}\\
\textbf{Aufgabenbereich}:\\
\LaTeX{}\\
\textbf{Betreuer}:\\
BEd. Otto Schuller
\mbox{}\\
\mbox{}\\
\mbox{}\\
\mbox{}\\
\mbox{}\\
\newpage